\section{Limitations and Conclusion}
\label{sec:limitations}

AutoTerm-SST is a live terminology-aware SST system, not a new foundation
model or topic classifier. The live demo serves En$\rightarrow$Zh/Ja/De;
routing stays correct across languages but switches more slowly without
language-tuned topic keywords (Appendix~\ref{app:multilingual}). Large open memories improve coverage but reduce
retrieval precision under the 10-candidate cap, so the system routes among
domain slices and keeps broad memory as a fallback pool.
The hybrid router is intentionally conservative and may stay on
the initial domain slice when evidence is ambiguous, as in our smoke talk. Our evaluations are single-run system probes; they characterize the
live system rather than establish corpus-level significance. Finally, the real backend requires high-end NVIDIA GPUs and external
model/retriever artifacts; mock mode supports UI and protocol testing without
them.

We presented AutoTerm-SST, a zero-setup terminology-memory workbench for
streaming speech translation. It keeps broad terminology resources offline,
selects active inventory slices from streaming evidence, ranks candidates into
a capped top-10 prompt interface, and exposes terms, glossary state, router
decisions, and latency in real time --- improving terminology
recall with zero setup.
